\documentclass[11pt]{article}
\usepackage[margin=1in]{geometry}
\usepackage{amsmath,amssymb,amsthm,mathtools}
\usepackage[hidelinks]{hyperref}

\newtheorem{theorem}{Theorem}
\newtheorem{proposition}{Proposition}
\newtheorem{lemma}{Lemma}
\newtheorem{remark}{Remark}

\newcommand{\R}{\mathbb R}
\newcommand{\one}{\mathbf 1}
\newcommand{\M}{\mathcal M}
\newcommand{\Om}{\Omega}

\title{Multi-Center \(A_1\) Mixtures at the Rubio de Francia Endpoint}
\author{Run fa\_banach\_001, agent\_lane\_16}
\date{2026-08-11}

\begin{document}
\maketitle

\section*{Source question and status}

For a pairwise disjoint family of intervals \(\Om\) in \(\R\), let
\[
 T^{\Om}f(x)=\left(\sum_{\omega\in\Om}|H_\omega f(x)|^2\right)^{1/2},
 \qquad
 \widehat{H_\omega f}=\one_\omega\widehat f.
\]
After Theorem 5.2 of Amenta--Lorist--Veraar, arXiv:1703.06044, the
strong estimate at the scalar endpoint \(p=q'\) is recorded as open.  Its
central \(q=2\) case asks whether
\begin{equation}\label{eq:conjecture}
 \|T^{\Om}f\|_{L^2(w)}\lesssim_w\|f\|_{L^2(w)}
 \quad\text{for every }w\in A_1,
\end{equation}
uniformly in \(\Om\).  Di Plinio--Fl\'orez-Amatriain--Parissis--Roncal,
arXiv:2308.01442, still call \eqref{eq:conjecture} the outstanding
trigonometric endpoint conjecture.  They prove it when \(w\) is even and
radially decreasing, as well as for the Walsh analogue.

The result below closes their radial class under arbitrary positive mixtures.
It produces multi-center and multiscale weights which need not be radially
decreasing about any point.

\section*{Proof intuition}

The key observation is elementary but powerful: after squaring, an
\(L^2(w)\) estimate is linear in \(w\).  If the endpoint estimate holds with
one common constant for each member of a family \((v_s)\), Tonelli's theorem
lets us integrate the estimates and proves the same bound for every positive
mixture \(w=\int v_s\).  Uniform \(A_1\) control is also preserved because
the maximal operator is sublinear.  Combining this cone principle with the
2023 radial theorem gives the claimed class.

\begin{lemma}[positive-mixture principle]\label{lem:mixture}
Let \((S,\nu)\) be a measure space and let \(v_s\) be nonnegative locally
integrable weights, measurable in \((s,x)\).  Suppose that, with a common
constant \(B\),
\[
 \|T^{\Om}f\|_{L^2(v_s)}\le B\|f\|_{L^2(v_s)}
\]
for \(\nu\)-almost every \(s\), every finite pairwise disjoint \(\Om\), and
every Schwartz function \(f\).  If
\(w(x)=\int_Sv_s(x)\,d\nu(s)\) is finite almost everywhere and locally
integrable, then
\[
 \|T^{\Om}f\|_{L^2(w)}\le B\|f\|_{L^2(w)}.
\]
If in addition \([v_s]_{A_1}\le K\) almost everywhere, then
\([w]_{A_1}\le K\).
\end{lemma}

\begin{proof}
Tonelli's theorem gives
\[
 \int_\R |T^{\Om}f|^2w
 =\int_S\int_\R |T^{\Om}f|^2v_s\,dx\,d\nu(s)
 \le B^2\int_S\int_\R |f|^2v_s\,dx\,d\nu(s)
 =B^2\int_\R|f|^2w.
\]
For the last assertion, positivity and sublinearity give, almost everywhere,
\[
 \M w\le\int_S\M v_s\,d\nu(s)
 \le K\int_Sv_s\,d\nu(s)=Kw.
\]
\end{proof}

\begin{theorem}[multi-center radial mixtures]\label{thm:main}
Fix \(K<\infty\).  For \(s\in S\), let
\[
 v_s(x)=\phi_s\!\left(\frac{|x-a_s|}{r_s}\right),
 \qquad a_s\in\R,\quad r_s>0,
\]
where \(\phi_s:[0,\infty)\to[0,\infty)\) is nonincreasing and
\([v_s]_{A_1}\le K\).  Let \(w=\int_Sv_s\,d\nu(s)\) be finite almost
everywhere and locally integrable.  There is \(C(K)<\infty\), independent of
the mixture and of the disjoint family \(\Om\), such that
\[
 \|T^{\Om}f\|_{L^2(w)}\le C(K)\|f\|_{L^2(w)}.
\]
Moreover \(w\in A_1\) and \([w]_{A_1}\le K\).
\end{theorem}

\begin{proof}
Translation and dilation preserve both the \(A_1\) characteristic and the
form of the Fourier-projection inequality.  Thus the radial endpoint theorem
of arXiv:2308.01442 applies to every \(v_s\), with a bound depending only on
\(K\) (its stated bound involves \([v_s]_{A_1}\) and
\([v_s]_{A_\infty}\), and the latter is controlled by the former).  Lemma
\ref{lem:mixture} finishes the proof.
\end{proof}

Here is a concrete, flexible subclass.

\begin{proposition}[polynomial-tail mixtures]\label{prop:kernels}
Fix \(0<\alpha<1\).  Let \(c_j\ge0\), \(a_j\in\R\), and \(r_j>0\), and
assume that
\[
 w(x)=\sum_{j}c_j
 \left(1+\frac{|x-a_j|}{r_j}\right)^{-\alpha}
\]
is finite almost everywhere and locally integrable.  Then \(w\in A_1\), with
\([w]_{A_1}\le K_\alpha\), and
\[
 \|T^{\Om}f\|_{L^2(w)}\le C_\alpha\|f\|_{L^2(w)}
\]
uniformly over all pairwise disjoint \(\Om\).  The same conclusion holds for
positive integral mixtures of these kernels.
\end{proposition}

\begin{proof}
It suffices to check that
\(k(x)=(1+|x|)^{-\alpha}\) has an \(A_1\) constant depending only on
\(\alpha\).  Fix an interval \(I\ni x\), put \(R=1+|x|\), and let
\(L=|I|\).  If \(L\le R/2\), then \(1+|y|\ge R/2\) for \(y\in I\), so
the average of \(k\) on \(I\) is at most \(2^\alpha k(x)\).  If
\(L>R/2\), then
\[
 \frac1L\int_I(1+|y|)^{-\alpha}\,dy
 \le C_\alpha L^{-\alpha}
 \le C'_\alpha R^{-\alpha}=C'_\alpha k(x).
\]
Translation and dilation do not change this bound.  The result follows from
Theorem \ref{thm:main}; integral mixtures follow from Tonelli.
\end{proof}

For example, two well-separated kernels in Proposition \ref{prop:kernels}
give two strict peaks, so the resulting weight is not radially decreasing
about any center.  Arbitrarily many centers and scales are allowed, and the
constant does not grow with their number.

\section*{A reduction toward the full conjecture}

The cone principle also narrows the remaining problem.  Write \(\M\) for the
uncentered Hardy--Littlewood maximal operator.

\begin{proposition}[Coifman--Rochberg reduction]\label{prop:reduction}
Fix \(K\).  There exist \(r>1\) and \(\delta\in(1/r,1)\), depending only on
\(K\), with the following property.  Every \(w\in A_1\) with
\([w]_{A_1}\le K\) is pointwise comparable, with constants depending only on
\(K\), to a positive mixture
\[
 v=\sum_{k\in\mathbb Z}2^k(\M\one_{E_k})^\delta,
 \qquad E_k=\{x:w(x)>2^k\}.
\]
Each \((\M\one_{E_k})^\delta\) has a common \(A_1\) bound.  Consequently,
to prove \eqref{eq:conjecture} for all weights of \(A_1\) characteristic at
most \(K\), it is enough to prove it uniformly for the special weights
\((\M\one_E)^\delta\), for arbitrary measurable \(E\subset\R\).
\end{proposition}

\begin{proof}
The reverse-H\"older theorem for \(A_1\) gives \(r>1\) and \(A=A(K)\) such
that
\[
 \frac1{|I|}\int_Iw^r\le A\left(\frac1{|I|}\int_Iw\right)^r.
\]
At almost every \(x\), for every interval \(I\ni x\), Chebyshev and the
\(A_1\) inequality imply
\[
 \frac{|E_k\cap I|}{|I|}
 \le 2^{-kr}\frac1{|I|}\int_Iw^r
 \le A K^r\left(\frac{w(x)}{2^k}\right)^r.
\]
Taking the supremum over \(I\) yields
\begin{equation}\label{eq:levelmax}
 \M\one_{E_k}(x)
 \le \min\left\{1,AK^r\left(\frac{w(x)}{2^k}\right)^r\right\}.
\end{equation}
Choose \(\delta\in(1/r,1)\).  At Lebesgue points in \(E_k\),
\(\M\one_{E_k}=1\).  Hence the terms with \(2^k<w(x)\) sum to a quantity
comparable from above and below to \(w(x)\).  For the remaining terms,
\eqref{eq:levelmax} and \(r\delta>1\) give a convergent geometric series:
\[
 \sum_{2^k\ge w(x)}2^k(\M\one_{E_k}(x))^\delta
 \le C(K)w(x)^{r\delta}
 \sum_{2^k\ge w(x)}2^{-k(r\delta-1)}
 \le C'(K)w(x).
\]
Thus \(w\lesssim v\lesssim_Kw\).  The standard Coifman--Rochberg lemma
states that \((\M g)^\delta\in A_1\), with constant depending only on
\(\delta\), for every nonnegative \(g\).  If the desired endpoint estimate
holds uniformly for these special weights, Lemma \ref{lem:mixture} proves it
for \(v\); pointwise comparability transfers it from \(v\) to \(w\).
\end{proof}

\begin{remark}[scope]
Theorem \ref{thm:main} is an unconditional partial result.  Proposition
\ref{prop:reduction} is a reduction, not a proof of the full conjecture: a
general weight \((\M\one_E)^\delta\) need not be a positive mixture of
single-center radial weights by any argument presently supplied here.  The
case of all \(A_1\) weights, and the source's other endpoint cases
\(p=q'<2\), remain open.
\end{remark}

\section*{Novelty check}

Exact-id, title, endpoint-phrase, \(p=q'\), \(A_1\), radial-weight, Walsh,
and positive-mixture searches were run in the local FA/Banach registry and
in primary arXiv sources through 2026-08-11.  The decisive later source was
arXiv:2308.01442.  No full endpoint resolution and no explicit statement of
the multi-center positive-mixture consequence above was found.  Since the
main proof is a short cone argument from the 2023 theorem, novelty confidence
is moderate while validity confidence is high.

\section*{References}

\begin{itemize}
\item A. Amenta, E. Lorist, and M. Veraar,
\emph{Rescaled extrapolation for vector-valued functions},
arXiv:1703.06044.
\item F. Di Plinio, M. Fl\'orez-Amatriain, I. Parissis, and L. Roncal,
\emph{Pointwise localization and sharp weighted bounds for Rubio de Francia
square functions}, arXiv:2308.01442.
\item R. R. Coifman and R. Rochberg,
\emph{Another characterization of BMO}, Proc. Amer. Math. Soc. 79 (1980),
249--254 (maximal-function powers in \(A_1\)).
\end{itemize}

\end{document}
